% 辛群（综述）
% license CCBY4
% type Wiki

本文根据 CC-BY-SA 协议转载翻译自维基百科\href{https://en.wikipedia.org/wiki/Symplectic_group}{相关文章}

在数学中，“辛群”的名称可以指两类不同但密切相关的数学群族，分别记作Sp(2n, F)与Sp(n)，其中n正整数，F为域（通常取C或 R）。后一种被称为紧致辛群，也记作$\mathrm{USp}(n)$。许多作者使用略有不同的记号，通常差别体现在系数 2 的处理上。此处采用的记号与表示这些群的最常见矩阵的规模保持一致。在卡当（Cartan）对单李代数的分类中，复群Sp(2n, C) 的李代数记作Cₙ，而Sp(n)则是Sp(2n, C)的紧致实型。需要注意的是，当我们提及“（紧致）辛群”时，通常是指依照维数n编号的一族（紧致）辛群。

“辛群”这一名称由赫尔曼·外尔提出，用来取代此前令人困惑的名称——“（线）复群”（(line) complex group）和“阿贝尔线性群”。它是“复数”一词的希腊语类比。

梅塔辛群是实数域上辛群的双覆盖；它在其他局部域、有限域以及阿德尔环上也有类似的对应物。
\subsection{Sp(2n, F)}
辛群是一个经典群，定义为作用在2n 维向量空间（定义在域$F$上）上的线性变换的集合，这些变换保持一个非退化的斜对称双线性型。这样的向量空间称为辛向量空间，而一个抽象辛向量空间$V$的辛群记作$\mathrm{Sp}(V)$。一旦为$V$选定基底，辛群就变为由矩阵乘法构成的$2n \times 2n$辛矩阵群，其元素属于域$F$。这个群记作$\mathrm{Sp}(2n, F)$或$\mathrm{Sp}(n, F)$。如果该双线性型由一个非奇异斜对称矩阵 $\Omega$ 表示，则有：
$$
\operatorname{Sp}(2n,F) = \{ M \in M_{2n \times 2n}(F) : M^{\mathrm{T}} \, \Omega \, M = \Omega \},~
$$
其中$M^{\mathrm{T}}$表示矩阵$M$的转置。通常，$\Omega$定义为：
$$
\Omega = 
\begin{pmatrix}
0 & I_n \\
- I_n & 0
\end{pmatrix},~
$$
其中$I_n$是 $n \times n$ 单位矩阵。在这种情况下，$\mathrm{Sp}(2n, F)$ 可以表示为如下的分块矩阵：
$$
\begin{pmatrix}
A & B \\
C & D
\end{pmatrix}, \quad A, B, C, D \in M_{n \times n}(F),~
$$
并满足以下三个方程：
$$
\begin{aligned}
- C^{\mathrm{T}} A + A^{\mathrm{T}} C &= 0, \\
- C^{\mathrm{T}} B + A^{\mathrm{T}} D &= I_n, \\
- D^{\mathrm{T}} B + B^{\mathrm{T}} D &= 0.
\end{aligned}~
$$
由于所有辛矩阵的行列式都等于 1，辛群是特殊线性群$\mathrm{SL}(2n, F)$的一个子群。当$n = 1$时，矩阵满足辛条件当且仅当其行列式为 1，因此：$\mathrm{Sp}(2, F) = \mathrm{SL}(2, F)$.而当$n > 1$时，还存在附加条件，即$\mathrm{Sp}(2n, F)$是$\mathrm{SL}(2n, F)$的一个真子群。

通常，域$F$取实数域$\mathbf{R}$或复数域$\mathbf{C}$。在这些情况下，$\mathrm{Sp}(2n, F)$分别是实或复的李群，其实或复维数为：$n(2n+1)$.这些群是连通的，但不是紧致的。

当域的特征不为 2 时，$\mathrm{Sp}(2n, F)$的中心由矩阵$I_{2n}$ 和 $-I_{2n}$ 构成。\(^\text{[1]}\)由于其中心是离散的，且模去中心的商群是单群，因此 $\mathrm{Sp}(2n, F)$ 被视为一个单李群。

对应李代数的实秩，因此李群$\mathrm{Sp}(2n, F)$的实秩为：$n$。

辛群$\mathrm{Sp}(2n, F)$的李代数定义为：
$$
\mathfrak{sp}(2n,F) = \{ X \in M_{2n \times 2n}(F) : \Omega X + X^{\mathrm{T}} \Omega = 0 \},~
$$
其李括号由对易子给出。\(^\text{[2]}\)

对于标准的斜对称双线性型：$\Omega = \begin{pmatrix}0 & I \\- I & 0
\end{pmatrix}$，该李代数由所有分块矩阵$\begin{pmatrix}A & B \\C & D
\end{pmatrix}$构成，并满足条件：
$$
\begin{aligned}
A &= - D^{\mathrm{T}}, \\
B &= B^{\mathrm{T}}, \\
C &= C^{\mathrm{T}}
\end{aligned}~
$$
\subsubsection{Sp(2n, C)}
复数域上的辛群是一个非紧致、单连通的单李群。该群的定义中并不涉及共轭（与直觉上的预期相反），而只是将定义中的域换为复数域，其余保持完全相同。\(^\text{[3]}\)
\subsubsection{Sp(2n, ℝ)}
$\mathrm{Sp}(n, \mathbf{C})$是实群$\mathrm{Sp}(2n, \mathbf{R})$ 的复化。$\mathrm{Sp}(2n, \mathbf{R})$是一个实的、非紧致的、连通的单李群。\(^\text{[4]}\)它的基本群同构于加法下的整数群。作为单李群的实型，它的李代数是可分解李代数。

$\mathrm{Sp}(2n, \mathbf{R})$的进一步性质：
\begin{itemize}
\item 从李代数$\mathfrak{sp}(2n, \mathbf{R})$ 到群 $\mathrm{Sp}(2n, \mathbb{R})$的指数映射不是满射。然而，该群的任意元素都可以表示为两个指数的乘积。\(^\text{[5]}\)换句话说：
   $$
   \forall S \in \operatorname{Sp}(2n,\mathbf{R}) \;\; \exists X, Y \in \mathfrak{sp}(2n,\mathbf{R}) \;\; S = e^{X} e^{Y}.~
   $$
\item 对于所有$S \in \mathrm{Sp}(2n, \mathbf{R})$，有：
   $$
   S = O Z O',~
   $$
   其中
   $$
   O, O' \in \operatorname{Sp}(2n, \mathbf{R}) \cap SO(2n) \cong U(n),~
   $$
   且
   $$
   Z = \begin{pmatrix} D & 0 \\ 0 & D^{-1} \end{pmatrix}.~
   $$
   这里 $D$ 是正定对角矩阵。所有这样的 $Z$ 构成 $\mathrm{Sp}(2n, \mathbf{R})$ 的一个非紧子群，而 $U(n)$ 构成一个紧子群。这种分解被称为欧拉分解或Bloch–Messiah 分解。\(^\text{[6]}\)更多关于辛矩阵的性质可参见维基百科相关页面。
\item 作为李群，$\mathrm{Sp}(2n, \mathbf{R})$具有流形结构。该流形与酉群$U(n)$和一个维数为$n(n+1)$的向量空间的笛卡尔积微分同胚。\(^\text{[7]}\)
\end{itemize}
\subsubsection{无穷小生成元}
辛李代数$\mathfrak{sp}(2n, F)$的元素称为哈密顿矩阵。

这些矩阵$Q$的一般形式为：
$$
Q = 
\begin{pmatrix}
A & B \\
C & -A^{\mathrm{T}}
\end{pmatrix},~
$$
其中$B$与$C$是对称矩阵。推导过程可参见“经典群”。
\subsubsection{辛矩阵的例子}
对于$\mathrm{Sp}(2, \mathbb{R})$，即所有行列式为 1 的$2 \times 2$ 矩阵，该群的三个辛$(0,1)$-矩阵为：\(^\text{[8]}\)
$$
\begin{pmatrix}
1 & 0 \\
0 & 1
\end{pmatrix}, \quad
\begin{pmatrix}
1 & 0 \\
1 & 1
\end{pmatrix}, \quad
\text{以及} \quad
\begin{pmatrix}
1 & 1 \\
0 & 1
\end{pmatrix}.~
$$
